% SAME-FAMILY WORKING FRAGMENT, 25 October 2026. Root-owned.
% Complete A17+ordinary A1 / actual I transcription; outside review pending.
% Unintegrated and uncompiled. No new bound or runtime credit.
\subsection{The two rank-one spectators beside \(A_{17}\)}
\label{app:A17rankone18}
We now exclude the two genuinely different rank-one alternatives.
An ordinary \(A_1\) has energy greater than \(.1\); an actual isolated
short \(I\) need not. The latter supplies a wall and additional actual
edge pressures instead. Throughout this subsection the entire incident
forest is \(A_{17}+A_1\), at the original farthest vertex \(p\).
Write \(z=E(p-x-ye)\), where \(e\) is its unit rank-one normal,
\(x\) belongs to zero-sum \(\mathbb R^{18}\), and
\(\norm x^2+y^2<1\). Set
\[
 \Omega_i=\sqrt2(P_i-iP_{18}/18),\quad
 D_i=i(18-i)/9,\quad C_{ij}=i(18-j)/9\quad(i\le j).
\]
There is no arithmetic hypothesis.

\subsubsection{Common geometry and the unscaled references}
Call \(5\le i\le13\) weak when
\(\ip x{\Omega_i}\le\sqrt{D_i}-.61\).
The target-independent inventory of Appendix~\ref{app:Acentralwindow}
leaves near-four roles at \(5,6\) or their reflections, and pure helpers
\[
 c=Ap-g\Omega_i-\Omega_a/2,\quad a=1\text{ or }17,
 \quad\phi/2\le g<17/16,
\]
with own five, finite at least seven, or dotted. All foreign and own-zero
roles have their dispositions in that inventory. Pure helpers are
nonadjacent before near-four caps are removed. The cutoff \(17/16\)
concerns these weak helpers only, never arbitrary endpoint caps.

We remove near-four caps without any contrary component-energy target.
For a left cap its same-wall mate is the actual edge endpoint at
\(j=4\) or \(5\). Its full positive finite projection retains
\(A_{j-1}+B_{18-j}+A_1\), with
\[
 \delta=1/(18-j),\quad
 \sigma=1/2+(17-j)/((18-j)\sqrt2),\quad
 C_{\rm mate}=\sigma\sqrt{(18-j)D_j}>10.
\]
Indeed \(\sigma>1.14\) and the last root is at least \(28/3\).
Farthest comparison gives height square greater than \(9D_j/11\),
so the heights at four and five exceed \(2.25\) and \(2.43\).
At the opposite end there is a mate or a strong twelve/thirteen:
otherwise two adjacent weak pure helpers would remain. Their heights
exceed \(2.218\) and \(2.077\). Disjoint end-segment Cauchy has
worst case
\[
 \norm x^2>\tfrac12(2.43^2/5+2.218^2/6)
                  =1.000450333\ldots>1.
\]
The other placements and two mates give larger bounds. Reflection
excludes the right near-four. Thus every remaining weak site is a
nonadjacent pure helper, at the correct original comparison factor one.

Use the unscaled \(N=18\) reference of
Appendix~\ref{app:Aendpointcores}, with
\[
 H_5=H_{13}=2.077,\ H_6=H_{12}=C=2.218,
 \ H_7=H_{11}=2.314,\ H_8=H_{10}=2.371,\ H_9=2.389.
\]
Flatten a weak end through six and delete every other weak knot by
its neighboring chord. Its positive strong-site coefficients give
\(\ip xw\ge E=\norm w^2\). The complete omission table is
\[
\begin{array}{c|c|c}
\text{weak ends }5,13&E\text{ lower bound}&\text{end blocks}\\\hline
0&E_0=.8934308&5/5\\
1&E_1>.86256&6/5\\
2&E_2=9978827/12000000&6/6.
\end{array}
\]
Each six-block anchor has capacity at least \(821/6000\), and each
five-block anchor at least \(343/2500\), after every internal omission.
Their signed coordinate means satisfy \(b_6<.2614\), \(b_5<.294\).

We use only the explicit comparison-one core and block algebra printed
in Appendices~\ref{app:Aendpointcores} and~\ref{app:B18localA17},
not the latter subsection's neighbouring-vertex conclusion or its
\(1.061\) comparison hypothesis. To specify this reuse, a weak-end
foot has \(U=v-k n_i-\ell n_a-m n_j\), with corrections
\[
 (k,\ell,m)=(g,2/3,1/3),\quad
 (g,2(1-g^2),0),\quad(1/g,0,0)
\]
for own five, finite at least seven and dotted, respectively.
The first uses its actual finite \(H_4\) projected point; the other
two are future negative cores. Their positive original-root
coefficients check every omitted wall, including \(e\) and its
opposite. Their pressures \(\ip xU\ge H\) have far time
\(A_0^2=(65/9)g^2+(5/9)g-19/36\); near adds \(8g/9\).
On two six-blocks the projections are
\[
 \Pi U=(5g/6)\Omega_6+\Omega_{12}/12\quad(\text{far}),\qquad
 \Pi U=(5g/6+1/12)\Omega_6\quad(\text{near}).
\]
The residual \(Y=U-\Pi U\) has zero sum separately on both end
blocks and \(\ip wY=0\). Its far squares are
\[
 Q_5=2g^2/3+1/12,\quad
 Q_f=4g^4-16g^2/3+29/12,\quad
 Q_d=5g^2/3-19/12+g^{-2};
\]
near adds \(g/3\). The full scalar seals proved there give the
six/six deficit \(\Delta=H-C(5g/6+1/12)>.3405\) for finite
owns and \(>.37\) for every \(1<g<17/16\).
For a far six/five foot its deficit decreases by \(343/15000\)
and its square by \(1/60\); near is unchanged.

For clarity, the finite-at-least-nine derivative in that seal is a
derivative of the lower surrogate
\[
 \overline\Delta=A_0-1/4-C(5g/6+1/12)-343/15000,
\]
not of the varying-radical deficit. It satisfies
\(\overline\Delta'>.839\), \(Q_f'\le16/3\).
At \(g=.93\), \(\overline\Delta>.32035\) and the adjusted
square is \(.77940804\). Thus a far finite foot at weight
\(.145\) gains more than \(.076\); five/seven/eight give
the respective deficit/square pairs
\((.33915,.505),(.32655,.74),(.3177248,.763)\), which give
the same strict bound directly. A far dotted foot at weight \(.13\)
gains more than
\[
 .26(5207/15000)-.0169(1751/1500)
           =1057901/15000000>.0705267.
\]

The full two-foot crosses are also coordinate identities. In one
six-block write
\[
 a=(g\sqrt2/6)(P_5-5e_6)-k n_5,\qquad
 b=(e_1-P_6/6)/\sqrt2-\ell n_1-m n_2.
\]
Then \(a^2=5g^2/3+k^2-2gk\),
\(b^2=5/12+\ell^2+m^2-\ell m-\ell\le5/12\), and
\(\ip{a_i}{b_j}=g_i/6\). Equal orientations have total square
\(Q_{f,L}+Q_{f,R}+(g_L+g_R)/3\); mixed orientations have
\(Q_{f,L}+Q_{f,R}+2\ip{b_L}{b_R}+2g_{\rm near}/3\).
Hence the equal/mixed bounds are \(17/6,11/3\) for two finite
owns and \(3.077,3.91\) if a dotted own occurs. Unequal
parameters and both mixed orientations are included.

Every use below pays \(\mu\Pi U\) from its individual anchors
before using the actual pressure. For \(W=w+\mu\sum Y\), one
completed square therefore gives
\[
 \norm x^2\ge E+2\mu\sum\Delta-\mu^2\norm{\textstyle\sum Y}^2.
\]
No independently optimized component estimates are added.

\subsubsection{An ordinary rank-one spectator}
Suppose first that \(\norm x^2\le.9\). With one weak end,
weight \(.1\) fits both anchor capacities and gives gain
\(>.2(.317)-.01(8/5)=.0474\); hence
\(\norm x^2>.90996\). With two weak ends, both individual
prices are at most \(101/960<821/6000\). The full mixed square
is less than \(32/5\), giving gain
\(>.4(.34)-.01(32/5)=.072\), and total \(>.90356\).
Only zero weak ends remain.

Both five-block means are \(b<.294\). If one signed endpoint
coordinate were at least \(.4\), replace that block by
\(.4,r,r,r,r\), with \(r=(5b-.4)/4>.2625\).
The next coordinate is less than \(.1\), so all surviving
coefficients remain nonnegative. The block-sum-preserving gain is
\((5/4)(.4-b)^2>.014045\), contradicting
\(E_0+.014045>.9\). Thus both signed endpoint coordinates
are less than \(.4\), and both actual endpoint cosines are less
than \(3/2\): otherwise their squared coordinate demands are
at least \(17/90>.16\). This temporary graft gain is not reused.

For any endpoint wall \(c=Ap-v_A-\beta e\), the actual reciprocal
shadow and distinct orthogonal-opposite lemma give \(A>\beta\),
including exclusion of the boundary shadow one. Thus
\(\norm{v_A}^2>1\). Finite classification after deleting old one
leaves isolation, a three at \(j=2,3,16,17\), or a four at
\(j=2,17\). For \(v_A=g\Omega_1+\delta\Omega_j\), their
zero-foreign cosines, in that order, are
\[
 (34g+16)/18,\quad(34g+15)/(2\sqrt{18}),\quad
 (34g+2)/(2\sqrt{18}),\quad(34g+1)/18,\quad
 (34g+16\sqrt2)/6,\quad(34g+\sqrt2)/6.
\]
Foreign cost only increases the actual cosine. Own-zero cases violate
\(\norm{v_A}^2>1\) or have the actual affine \(E_8\) null at
three-at-three. Own three at two has its triangle null. The remaining
native substitutions and \(C<3/2\) leave only pure far-three at
seventeen, with own four. Indeed its \(1/\sqrt2\le g<13/17\)
allows only four; a foreign \(\beta^2\ge1/4\) multiplies its
cosine square by at least \(36/19\), giving \(>64/19>9/4\).
Two such opposite pure caps are distinct and have product
\(8\sqrt2/9-1/3>0\). At least one endpoint is therefore isolated
within \(A_{17}\), with foreign contact still allowed.

Put \(D=17/9\), \(Q=\beta^2\), \(u=(1-Q)/g^2\).
The inequality \(Dg^2>1\) removes own zero, three and four;
six and parallel have actual feasible nulls. Also \(C<3/2\)
gives \(u>68/81>.83\). The following complete point constructions
give a pure \(A\) displacement \(Z=\zeta-P_{18}/18\),
\(\sum\zeta=1\), \(\ip xZ\ge t_0\), \(R=\norm\zeta^2\).

For pure five use the finite \(H_4\) point from
\(c+(3g/2)n_1+gn_2+(g/2)n_3\), with
\(\delta=1-3g^2/2>0\). Here \(\zeta=(1/4,1/4,1/4,1/4)\),
\(R=1/4\), and
\(t_0=AL/(\sqrt2g(L+\sqrt\delta))>.33\), where
\(L^2=7g^2/18\), \(A>.486,L>.504,\sqrt\delta<.136\).
For pure finite at least seven use
\(c+g^{-1}n_1+2(g^{-1}-g)n_2\). Its selected products are
zero, zero, and \((3-4g^2)/(2g)<0\), and
\[
 \zeta=(1-u/2,1-u/2,u-1),\quad R=3-4u+3u^2/2,
 \quad t_0={\sqrt{D-u}-\sqrt{7u-4-3u^2}\over\sqrt2}.
\]
Seven has \(1.23<u<1.233\), \(R<.35,t_0>.38\).
Finite at least eight has \(1<u<1.1716\), \(R\le.5,t_0>.39\);
the two radicals are respectively \(>.8469\) and \(<.28868\).

For any \(Q\) with \(0<u<1\), clear the entire foreign contact
in \(c+\beta e+((1-Q)/g)n_1\). Its square is
\(-g^2u(1-u)\), its cap/foreign products are zero and its own
product is negative. Its positive time and original coefficients
check every other wall. It gives
\[
 \zeta=(1-u/2,u/2),\quad R=1-u+u^2/2,\quad
 t_0={\sqrt{D-u}-\sqrt{u(1-u)}\over\sqrt2}>.4,
 \quad R<.515.
\]
The last bounds use \(u>.83\), radicals \(>.9428\), \(<.376\).
The boundary \(u=1\) is an actual future null.
For \(Q>0,u>1\), the native alphabet forces own five,
foreign three and \(Q=1/4\). Its positive finite correction
\(c+gn_1+e/2\) has square \(\delta=3/4-g^2>0\).
Its actual projected point gives \(R=1/2\) and \(t_0>.42\),
using \(A>.697,L>.762,\sqrt\delta<.31\),
\(L^2=8g^2/9\). These are points, not native projected normals.

At the isolated end project \(Z\) onto \(\Omega_5/(5\sqrt2)\).
Weight \(1/2\) costs less than \(.071<343/2500\), and the
residual has square \(R-1/5\), orthogonal to the reference.
Its gain is greater than \(t_0-.294-(R-1/5)/4\).
For the five rows just listed these strict gains exceed
\(.0235,.0485,.021,.02725,.051\). Each is greater than
\(.9-E_0=.0065692\). This closes the zero-weak-end branch.
We have proved \(\norm x^2>.9\) for the entire rank-one forest.
An ordinary spectator has \(y^2>.1\), giving the contradiction.

\subsubsection{The actual isolated short spectator}
Now let its actual wall be
\(c_I=\sqrt h\,p-\sqrt{1+h}\,e\), \(0<h<1/2\).
Appendix~\ref{app:Ijoint} proves, on this same original vector,
\[
 y\ge s,\quad x_1,-x_{18}\ge a_I=\sqrt{17/18}\,t,
 \quad t+\tfrac32s>.767.
\]
Put \(k=\sqrt{17/18}\), \(D_0=.7453706\). Since \(k>.9718\),
\(a_I+(3k/2)s>D_0\). Weighted Cauchy therefore gives, for
\(c>0,B<D_0\),
\begin{equation}\label{eq:A17Igraft}
 s^2+c(a_I-B)_+^2>{(D_0-B)^2\over17/8+1/c}.
\end{equation}
It counts the spectator energy once. An endpoint graft on a block
of length \(L\), from its current coordinate \(B\) toward \(a_I\),
gains \(L(a_I-B)_+^2/(L-1)\) and costs
\((a_I-B)_+/((L-1)\sqrt2)\) at that individual anchor.
It may be capped at a smaller coordinate when necessary. Its full
cross with previously paid helper residuals is included in \(B\).

For two weak ends pay \(.1\) to each foot. Equal orientations
spend at most \(31/320\) at either anchor and have both endpoint
coordinates at most \(.3454\). Mixed orientations spend at most
\(101/960\) and \(17/192\); the unloaded coordinate is at most
\(.2865\). If \(a_I\le.6\), the respective additional graft
prices are at most \(.2546/(5\sqrt2)\) and
\(.3135/(5\sqrt2)\). These fit their separate remaining capacities
\(821/6000-31/320\) and \(821/6000-17/192\).
Equation~\eqref{eq:A17Igraft} gives gains \(>.0629\) for
\(c=12/5,B=.3454\), and \(>.0711\) for
\(c=6/5,B=.2865\). If \(a_I>.6\), cap those grafts at
\(.6\); their gains alone exceed \(.155,.117\).
For two finite owns the base totals exceed
\(E_2+.1362-.01(17/6)>.9394355\) and
\(E_2+.1362-.01(11/3)>.9311022\). A dotted own raises the
deficit sum above \(.7105\) and gives larger totals. The appropriate
graft therefore makes each full total greater than one.

For one weak end orient it to five. A near finite helper at weight
\(.145\) spends less than \(.145(11/12)<821/6000\), has square
\(<17/12\) and gain \(>.06895\). A near dotted helper at
\(.14\) spends less than \(.14(31/32)\), has square \(<1.539\)
and gain \(>.0734\). In each case the other five-block is unloaded.
Its full graft is legal for \(a_I<.972\), since
\((.972-.294)/(4\sqrt2)<.12<343/2500\).
Equation~\eqref{eq:A17Igraft}, \(c=5/4,B=.294\), gives
graft plus spectator \(>.0696\). Thus even the smaller total
is \(>.86256+.06895+.0696>1.0011\).

A far finite helper at weight \(.145\) has gain \(>.076\).
Its own price leaves more than \(.016\), and its opposite price
leaves at least \(.1227\),
and its endpoint coordinates are at most \(.2956,.3761\).
The weak graft is legal up to increment \(.113\); the other
graft is legal throughout \(a_I<.972\). If \(a_I\le.4086\),
the weak graft and \eqref{eq:A17Igraft} give \(>.0683\),
with \(c=6/5,B=.2956\). Above that split cap the weak gain
at \((6/5)(.113)^2=.0153228\), and use the strong graft
with spectator \(>.0466\), \(c=5/4,B=.3761\).
Both give an extra \(>.0619\) and total \(>1.00046\).

A far dotted helper at weight \(.13\) has gain \(>.0705267\),
leaves at least \(1043/48000\) and \(.1242\) at the two anchors,
and has endpoint coordinates at most \(.294,.3676\).
The weak graft is legal up to increment \(.1535\).
Below \(a_I=.4475\), its graft plus spectator is \(>.0688\),
with \(c=6/5,B=.294\). Above this split its capped gain is
\((6/5)(.1535)^2=.0282747\); the opposite graft with spectator
is \(>.0487\), with \(c=5/4,B=.3676\).
Both totals exceed one. Reflection covers the other weak end.

Finally suppose neither end is weak. One full five-block graft with
the spectator gives \(>.0696\). If the other actual endpoint
cosine were at least \(8/5\), its signed coordinate would exceed
\(7/15\), since \(17/78>(7/15)^2\). The disjoint grafts give
\[
 \norm x^2+y^2>E_0+.0696+(5/4)(7/15-.294)^2>1.000298.
\]
Thus both endpoint cosines are less than \(8/5\). The complete
foreign-safe list already proved above leaves only isolation or pure
far-three with own four/five: now \(g<139/170\), and a foreign
contact would still give \(C^2>64/19>(8/5)^2\).
Two such far caps are distinct; their spatial product exceeds
\(17(g_L+g_R)/18>1.32\), while each time square
\((68g^2+4g-19)/36<.81\). Their actual product is positive.
At least one endpoint is isolated within \(A_{17}\).

For that isolated endpoint pure five is impossible: its time square
\((17\phi-19)/36<.237\) makes its strictly negative product with
\(c_I\) smaller in magnitude than \(1/2\). The foreign
five/three row with \(u>1\) is also impossible, since acuteness
would require \(h(Dg^2-1)\ge1/4\), whereas
\(h<1/2,Dg^2-1<.237\). All other rows use the same actual cores
as above, without the ordinary branch's temporary endpoint bound.
Pure seven gives gain \(>.0485\). For finite at least eight,
split \(1<u<1.1716\) at \(1.1\): the bounds
\((R,t_0)\le(.5,\text{ lower } .4407)\) and
\((.415,\text{ lower } .3945)\) give gains \(>.0717,.04675\).

For every \(0<u<1\) the cleared-core gain is \(>.06\).
Here are complete continuous seals, so no sampled parameter is used.
Lower chords for \(\sqrt{D-u}\) are
\(1.3743-.3916u,1.4012-.4454u,1.4402-.4974u\)
on \([0,1/2],[1/2,3/4],[3/4,1]\); positive endpoint squares
prove them. Upper chords for \(R\) are
\(1-.75u,.8125-.375u,.625-.125u\).
With \(.7071<1/\sqrt2<.7072\), the gain is bounded below by
\[
\begin{array}{c|c}
[0,1/2]&.47776753-.0895u-.7072\sqrt{u(1-u)}\\
[1/2,3/4]&.54366352-.2212u-.7072\sqrt{u(1-u)}\\
[3/4,1]&.61811542-.3205u-.7072\sqrt{u(1-u)}.
\end{array}
\]
The inequality \(au+b\sqrt{u(1-u)}\le(a+\sqrt{a^2+b^2})/2\),
with radical upper bounds \(.713,.7411,.7765\), gives respectively
\(>.0765,.0625,.0696\). Thus every surviving isolated endpoint
supplies gain \(>.0467\), paid at its own anchor by weight \(1/2\).
At the opposite end the independent full graft plus the one spectator
energy gives \(>.0696\). Their residuals have disjoint support.
One square yields
\[
 \norm x^2+y^2>E_0+.0467+.0696=1.0097308>1.
\]
This excludes the actual isolated short spectator as well. All original
foreign contacts, own parameters and wall signs have been retained;
the same original farthest vector and separate anchor prices are used
throughout. Both rank-one alternatives beside \(A_{17}\) are now closed.
